% TOP_PROBLEM: {"id": 129, "title": "Cubic Curves in F_p^2", "category_id": 6, "category": {"id": 6, "name": "geometry", "display_name": "Geometry", "description": "Euclidean and non-Euclidean geometry, geometric structures.", "slug": "geometry", "order_index": 6, "created_at": "2026-07-31T15:26:25.671Z"}, "status": "open"}
% FIRST_POSTED: null
% ENTRY_KIND: statement_only
% Imported from: batch03.tex; UnsolvedMath ID 129
\documentclass[11pt]{article}
\usepackage[margin=25mm]{geometry}
\usepackage{fontspec}
\setmainfont{Latin Modern Roman}
\usepackage{amsmath,amssymb,mathtools}
\usepackage{unicode-math}
\setmathfont{Latin Modern Math}
\usepackage{xurl,hyperref}
\hypersetup{colorlinks=true,urlcolor=blue}
\setcounter{secnumdepth}{0}
\setlength{\parindent}{0pt}
\setlength{\parskip}{5pt}
\setlength{\emergencystretch}{3em}
\newcommand{\sref}[2]{\hyperlink{source:#1:#2}{[S#2]}}
\newcommand{\cataloguescope}[1]{\paragraph{Recorded target.}#1}
\begin{document}
% UnsolvedMath ID 129; source catalogue code GREEN-041
\setcounter{section}{128}
\section{Cubic Curves in $\mathbb F_p^2$}\label{129}
{\small\sffamily UnsolvedMath ID 129 · Geometry}
\subsection{Definitions and mathematical statement}

\paragraph{Shared notation.}$\mathbb N=\{1,2,\ldots\}$, and $\mathbb Z,\mathbb Q,\mathbb R,\mathbb C$ denote the integers, rationals, reals and complex numbers. Any other conventions are specified locally.


For a prime $p$, let $\mathbb F_p$ be the field with $p$ elements. An affine line in $\mathbb F_p^2$ is a set $\{v+t w:t\in\mathbb F_p\}$ with $w\ne0$. A cubic curve here is the zero set of a nonzero polynomial $F\in\mathbb F_p[X,Y]$ of total degree at most three; reducible polynomials are permitted. A curve of smaller positive degree can also be included in a degree-three zero set by multiplying its defining polynomial by linear factors.
\paragraph{Statement recorded in the supplied source.}
For every real $\epsilon>0$, is there $p_0(\epsilon)$ such that for every prime $p\geq p_0(\epsilon)$ and every set $A\subset\mathbb F_p^2$ satisfying $|A\cap\ell|\leq2$ for every affine line $\ell$, some nonzero polynomial $F\in\mathbb F_p[X,Y]$ of total degree at most three satisfies
\[
 |\{a\in A:F(a)\ne0\}|\leq\epsilon p?
\]
Thus the exceptional $o(p)$ bound must be uniform over the sets $A$, as $p$ tends to infinity through primes. \sref{129}{2}
\subsection{Short English statement}
Does every planar set over a large prime field with no three points on a line lie almost entirely on a curve of degree at most three?
\subsection{Sources}
\begin{description}
\item[\hypertarget{source:129:1}{S1}] ULAM.AI and the UnsolvedMath Contributors, \emph{UnsolvedMath: A Curated Collection of Open Mathematics Problems}, record 129 (GREEN-041). CC BY 4.0. \url{https://huggingface.co/datasets/ulamai/UnsolvedMath}.
\item[\hypertarget{source:129:2}{S2}] Ben Green, \emph{100 Open Problems}, maintained author collection, Problem 68 and its comments, accessed 7 October 2026. \url{https://people.maths.ox.ac.uk/greenbj/papers/open-problems.pdf}.
\end{description}
\label{end:129}
\end{document}
